\documentclass[11pt]{article}
\usepackage{../../latex/stat4ml_style}

\title{Chapter 9: Gaussian Processes and the Infinite-Width Limit \\ \large Part III --- Statistical Foundations of Deep Learning \& Reinforcement Learning}
\author{Omid Safarzadeh}
\date{}

\begin{document}
\maketitle

\begin{abstract}
\noindent
A Gaussian process is a probability distribution not over numbers but over
\emph{functions}: any finite set of its values is jointly Gaussian, with
covariance set by a kernel that encodes which inputs should have similar outputs.
This makes nonparametric regression exact and Bayesian --- condition on data and
the posterior is again a GP, giving not just a predicted curve but a calibrated
uncertainty that is small near the data and large away from it. The chapter
closes with the result that ties this back to deep learning: a neural network of
infinite width, with random weights, \emph{is} a Gaussian process. The central
limit theorem does it --- each output is a sum of many independent random
contributions --- so the Bayesian prior implied by a wide network is a GP with a
kernel you can write down.
\end{abstract}

\tableofcontents

% ======================================================================
\section{A distribution over functions}
% ======================================================================

\begin{definition}[Gaussian process]
A \keyterm{Gaussian process} $f\sim\mathcal{GP}(m,k)$ is a collection of random
variables $\{f(x)\}$ such that any finite subset $(f(x_1),\dots,f(x_n))$ is
jointly Gaussian, with mean $m(x_i)$ and covariance
$\operatorname{Cov}(f(x_i),f(x_j))=k(x_i,x_j)$.
\end{definition}

The \keyterm{kernel} $k$ is the whole model. The common choice, the \keyterm{RBF}
(squared-exponential) kernel $k(x,x')=\exp(-\|x-x'\|^2/2\ell^2)$, says nearby
inputs are strongly correlated and distant ones nearly independent, so samples
from the prior are smooth --- the left panel draws five of them. A GP is
nonparametric: it does not fix a functional form with finitely many parameters,
it places a prior directly on the space of functions through the kernel.

% ======================================================================
\section{Regression is conditioning}
% ======================================================================

Because everything is jointly Gaussian, learning is just conditioning a
multivariate normal (Part~I, Chapter~4). Given training inputs $X$ with noisy
outputs $y$ and a test point $x_\star$, stack them; the joint is Gaussian, and the
conditional $f(x_\star)\mid y$ is Gaussian with closed-form mean and variance.

\begin{theorem}[GP predictive distribution]
With kernel matrix $K=k(X,X)$, noise $\sigma^2$, and cross-covariance
$k_\star = k(x_\star,X)$,
\[
  \mathbb{E}[f(x_\star)\mid y] = k_\star (K+\sigma^2 I)^{-1} y,
  \qquad
  \operatorname{Var}[f(x_\star)\mid y] = k(x_\star,x_\star) - k_\star (K+\sigma^2 I)^{-1} k_\star^\top.
\]
\end{theorem}
\begin{proof}
$(f(x_\star), y)$ is jointly Gaussian with covariance blocks
$k(x_\star,x_\star)$, $k_\star$, and $K+\sigma^2 I$. The Gaussian conditioning
formula --- for jointly normal $(a,b)$, $\mathbb{E}[a\mid b]=\mu_a +
\Sigma_{ab}\Sigma_{bb}^{-1}(b-\mu_b)$ and
$\operatorname{Var}[a\mid b]=\Sigma_{aa}-\Sigma_{ab}\Sigma_{bb}^{-1}\Sigma_{ba}$
--- applied with $\mu=0$ gives exactly the two expressions.
\end{proof}

The middle panel shows the result: the posterior mean interpolates the data, and
the $\pm 2$-standard-deviation band pinches to nearly zero at each observation and
balloons where there are none. This calibrated ``I don't know away from the
data'' is what a plain network's point prediction lacks --- and it comes for free
from conditioning a Gaussian. The cost is the $(K+\sigma^2 I)^{-1}$, an $n\times
n$ solve, which is why exact GPs are reserved for smaller $n$.

% ======================================================================
\section{A wide network is a Gaussian process}
% ======================================================================

The link to deep learning is a theorem of Neal (1996). Take a one-hidden-layer
network with $H$ units,
\[
  f(x) = \frac{1}{\sqrt H}\sum_{j=1}^H v_j\,\phi(w_j^\top x + b_j),
\]
and put independent zero-mean priors on the weights $v_j,w_j,b_j$.

\begin{theorem}[Infinite-width networks are GPs]
As $H\to\infty$, the function $f$ converges to a Gaussian process. For any inputs,
the outputs $(f(x_1),\dots,f(x_n))$ are jointly Gaussian with covariance
$k(x,x')=\mathbb{E}\big[v^2\big]\,\mathbb{E}_{w,b}\big[\phi(w^\top x+b)\phi(w^\top x'+b)\big]$.
\end{theorem}
\begin{proof}[Proof sketch]
For fixed inputs, $f(x)$ is a sum of $H$ i.i.d.\ terms $v_j\phi(w_j^\top x+b_j)$,
each with mean zero and finite variance, scaled by $1/\sqrt H$. By the central
limit theorem (Part~I, Chapter~5) the sum converges to a Gaussian; applied jointly
to several inputs (the multivariate CLT) the whole vector is jointly Gaussian, so
$f$ is a GP. Its covariance is the per-unit expectation above, since the units are
independent.
\end{proof}

So the central limit theorem of Part~I resurfaces as a statement about neural
networks: the prior over functions implied by a wide, randomly-initialised network
is a Gaussian process with a specific kernel --- the \keyterm{neural network
Gaussian process} (NNGP) kernel. The right panel is the theorem happening: the
distribution of a single output over random weights is sharply non-Gaussian at
width $1$, and by width $200$ it has become the Gaussian the CLT predicts. This is
the bridge between the two halves of the course --- the Bayesian, kernel view of a
neural network is literally a Gaussian process --- and it is why GPs, kernels, and
deep networks are studied together.

\begin{summarybox}
A Gaussian process is a distribution over functions: any finite set of values is
jointly Gaussian with covariance given by a kernel $k$, which encodes smoothness
(e.g.\ the RBF kernel). Regression is Gaussian conditioning (Chapter~4): the
predictive mean $k_\star(K+\sigma^2 I)^{-1}y$ interpolates the data and the
predictive variance gives calibrated uncertainty, small at the data and large
away from it, at the cost of an $n\times n$ solve. The tie to deep learning is
Neal's theorem: a one-hidden-layer network with random weights and
$1/\sqrt H$ scaling converges, as $H\to\infty$, to a Gaussian process by the
central limit theorem (Chapter~5), with the NNGP kernel as its covariance --- the
Bayesian prior of a wide network is a GP.
\end{summarybox}

\end{document}
